% =========================================================================
% 국문 논문 조판 프리앰블 — LuaLaTeX + luatexko 전용.
%
% 빌드:  lualatex main_ko  →  bibtex main_ko  →  lualatex main_ko ×2
%        (목차 쪽번호가 맞으려면 마지막 패스가 한 번 더 필요하다.)
%
% 본문 라틴은 Latin Modern Roman, 한글은 함초롬바탕(HCR Batang).
% 이전 판은 XeLaTeX + xeCJK + 맑은고딕이었다. 고딕 본문은 발표자료로 읽히고
% 논문으로는 읽히지 않아, 본문 서체를 바탕(명조) 계열로 바꾸었다.
% =========================================================================

\usepackage[a4paper,top=25mm,bottom=25mm,left=27mm,right=27mm,footskip=12mm]{geometry}
\usepackage{luatexko}
\usepackage{amsmath,amssymb,bm}
\usepackage{microtype}
\usepackage{graphicx}
\usepackage{booktabs}
\usepackage{multirow}
\usepackage{enumitem}
\usepackage{xcolor}
\usepackage{tabularray}
\UseTblrLibrary{booktabs}
\usepackage{tikz}
\usetikzlibrary{arrows.meta,positioning}
\usepackage[font=small,labelfont=bf,labelsep=period]{caption}
\usepackage[numbers,sort&compress]{natbib}
\usepackage[hidelinks,unicode]{hyperref}

% ── 폰트 ─────────────────────────────────────────────────────────────────
% 화살표·수학기호·박스문자는 한글 폰트에 없다. luaotfload 대체 목록으로 메운다.
\directlua{
  luaotfload.add_fallback("symfallback", {
    "HCR Batang:mode=node;",
    "DejaVu Sans:mode=node;",
    "Segoe UI Symbol:mode=node;",
    "Malgun Gothic:mode=node;",
  })
}
\setmainfont{Latin Modern Roman}[
  RawFeature={fallback=symfallback},
  Numbers=Lining,
]
\setsansfont{Latin Modern Sans}[RawFeature={fallback=symfallback}]
\setmonofont{Consolas}[Scale=0.92]
\setmainhangulfont{HCR Batang}[BoldFont={HCR Batang Bold},AutoFakeSlant=0.2,Scale=0.96]
\setsanshangulfont{HCR Dotum}[BoldFont={HCR Dotum Bold},AutoFakeSlant=0.2,Scale=0.96]
% 고정폭 한글은 라틴 고정폭의 정확히 두 배 너비여야 표의 열이 맞는다.
\setmonohangulfont{GulimChe}[AutoFakeBold=1.5,AutoFakeSlant=0.2,Scale=1.0116]

% ── 본문 조판 ────────────────────────────────────────────────────────────
\linespread{1.20}
\setlength{\parindent}{10pt}
\setlength{\emergencystretch}{2em}
% 참고문헌의 영문 표제는 하이픈 없이는 판면을 넘는다 — 하이픈을 막지 않는다.
\brokenpenalty=5000
\widowpenalty=10000
\clubpenalty=10000
% 가운뎃점은 이 원고에서 항목 구분자다 — 그 뒤에서 줄을 바꿀 수 있어야 한다.
\registerbreakableafter{"00B7}

\setlist{topsep=3pt,itemsep=1.5pt,parsep=0pt,leftmargin=*}
\setlist[itemize]{label=\textendash}

% 절 제목 — article 기본 서식을 유지하되 간격만 논문 비례로 좁힌다.
\makeatletter
\renewcommand\section{\@startsection{section}{1}{\z@}%
  {-3.0ex \@plus -1ex \@minus -.2ex}{1.8ex \@plus.2ex}%
  {\normalfont\large\bfseries}}
\renewcommand\subsection{\@startsection{subsection}{2}{\z@}%
  {-2.6ex\@plus -1ex \@minus -.2ex}{1.2ex \@plus .2ex}%
  {\normalfont\normalsize\bfseries}}
\renewcommand\subsubsection{\@startsection{subsubsection}{3}{\z@}%
  {-2.2ex\@plus -1ex \@minus -.2ex}{1.0ex \@plus .2ex}%
  {\normalfont\normalsize\itshape}}
\makeatother

% 인용 단락 — 논문의 중심 질문과 각 장의 결론 요약이 여기 들어간다.
\renewenvironment{quote}
  {\list{}{\leftmargin=1.6em\rightmargin=1.6em\topsep=4pt\parsep=2pt}%
   \item\relax\small}
  {\endlist}

% ── 구조화 초록·PICO 형 표제어 목록 ───────────────────────────────────────
% 표제어와 서술이 짝을 이루는 덩어리는 표가 아니라 목록으로 조판한다.
\newenvironment{fielddesc}[1][6.6em]
  {\begin{list}{}{%
      \setlength{\labelwidth}{\dimexpr#1-0.8em\relax}%
      \setlength{\leftmargin}{#1}%
      \setlength{\labelsep}{0.8em}%
      \setlength{\itemsep}{3pt}%
      \setlength{\parsep}{0pt}%
      \setlength{\topsep}{4pt}%
      \renewcommand{\makelabel}[1]{\hfill\bfseries ##1}}}
  {\end{list}}

% ── 표 기본형 ────────────────────────────────────────────────────────────
% booktabs 규칙, 세로줄 없음, 행 간격만 확보한다.
\SetTblrStyle{foot}{font=\footnotesize}
\SetTblrInner{rowsep=1.8pt,colsep=4pt,stretch=1.0}
\SetTblrInner[longtblr]{rowsep=1.8pt,colsep=4pt,stretch=1.0}
\SetTblrStyle{caption}{font=\small}
\SetTblrStyle{contfoot}{font=\footnotesize\itshape,halign=r}
\DefTblrTemplate{contfoot-text}{default}{(다음 쪽에 이어짐)}
\DefTblrTemplate{conthead-text}{default}{}
\DefTblrTemplate{capcont}{default}{}

% ── 논증 사슬 도해 스타일 ────────────────────────────────────────────────
% 화살촉 지정 Latex[...] 의 대괄호가 \begin{tikzpicture}[...] 의 선택인자
% 안에 있으면 첫 ']' 에서 인자가 끊긴다 — 그래서 여기서 미리 정의한다.
\tikzset{
  argstep/.style={draw=black!40, rounded corners=1.5pt,
                  inner xsep=2.4mm, inner ysep=2.0mm,
                  text width=136mm, align=left, font=\footnotesize},
  argkey/.style={argstep, draw=black, line width=0.7pt},
  argflow/.style={-{Latex[length=2.2mm,width=1.7mm]}, black!55, line width=0.6pt},
}

% ── article 클래스가 영어로 박아 둔 구조 표제 ─────────────────────────────
\renewcommand{\abstractname}{초록}
\renewcommand{\refname}{참고문헌}
\renewcommand{\figurename}{그림}
\renewcommand{\tablename}{표}
\renewcommand{\contentsname}{목차}
\renewcommand{\appendixname}{부록}
\renewcommand{\today}{{\number\year}년 {\number\month}월 {\number\day}일}

% ── 원고 약칭 ────────────────────────────────────────────────────────────
\newcommand{\todo}[1]{\textcolor{red}{[TODO: #1]}}
\newcommand{\Ti}{\ensuremath{T_i}}
\newcommand{\Vrot}{\ensuremath{V_\mathrm{rot}}}
\newcommand{\skillp}{\ensuremath{\mathrm{skill}_{\mathrm{PCHIP}}}}
\newcommand{\seqv}{\texttt{seq\_v2}}
\newcommand{\bthree}{\texttt{b3k8}}

\graphicspath{{figures/}}
